\begin{longtable}{ |p{6cm}|p{8cm}|}
\caption{\textit{Use Case Scenario} Registrasi Akun}\label{tab:use-case-scenario-registrasi-akun} 
\\
\hline
\textbf{\textit{Use Case Name}}
&Registrasi Akun 
\\ \hline

\textbf{\textit{Actor}}
&\textit{Staff Infrastructure}, \textit{Manager Infrastructure} 
\\ \hline

\textbf{\textit{Brief Description}}
&\textit{Actor} yang tidak memiliki akun untuk \textit{login} dapat melakukan registrasi akun
\\ \hline

\textbf{\textit{Entry Conditions}}
&\textit{Actor} telah mengakses menu registrasi akun
\\ \hline

\textbf{\textit{Flow of event}}
&1. \textit{Actor} membuka menu registasi akun
\newline
2. \textit{Actor} mengisi seluruh \textit{form input} regitrasi akun dan menekan tombol \textit{submit}
\newline
3. Sistem menampilkan \textit{modal} verifikasi \textit{email} dan mengirimkan kode OTP ke \textit{email} \textit{actor}
\newline
4. \textit{Actor} mengisi kode OTP yang didapatkan melalui \textit{email}
\newline
5. \textit{Actor} menunggu akun diaktifkan oleh \textit{Manager}
\\ \hline

\textbf{\textit{Alternative scenario}}
&\textit{Actor} mengisi \textit{input email} pada \textit{form registrasi} dengan \textit{email} selain domain perusahaan dan sudah terdaftar maka : Sistem menampilkan pesan error
\\ \hline

\textbf{\textit{Exit Conditions}}
&Sistem menampilkan menu \textit{login}
\\ \hline
\end{longtable}

%%%%%%%%%%%%%%%%%%%%%%%%%%%%%%%%%%%%%%%%%%%%%%%%
\begin{longtable}{ |p{6cm}|p{8cm}|}
\caption{\textit{Use Case Scenario} Ganti Password}\label{tab:use-case-scenario-ganti-password} 
\\
\hline
\textbf{\textit{Use Case Name}}
&Ganti Password
\\ \hline

\textbf{\textit{Actor}}
&\textit{Staff Infrastructure}, \textit{Manager Infrastructure}, \textit{Manager Human Resources}
\\ \hline

\textbf{\textit{Brief Description}}
&\textit{Actor} yang ingin mengganti password akun
\\ \hline

\textbf{\textit{Entry Conditions}}
&\textit{Actor} telah melakukan \textit{login}
\\ \hline

\textbf{\textit{Flow of event}}
&1. \textit{Actor} membuka menu profil dan menekan tombol \textit{change password}
\newline
2. \textit{Actor} mengisi seluruh \textit{form input} pergantian password dan menekan tombol \textit{submit}
\newline
\\ \hline

\textbf{\textit{Alternative scenario}}
&\textit{Actor} yang salah mengisi password lama maka sistem akan mengembalikan pesan kesalahan
\\ \hline

\textbf{\textit{Exit Conditions}}
&Sistem selesai memproses pergantian passowrd dan menampilkan pesan berhasil ganti password
\\ \hline
\end{longtable}

%%%%%%%%%%%%%%%%%%%%%%%%%%%%%%%%%%%%%%%%%%%%%%%%%%%%%%%%%%%%%%%%%
\FloatBarrier
\begin{longtable}{ |p{6cm}|p{8cm}|}
\caption{\textit{Use Case Scenario Login}}\label{tab:use-case-scenario-login} 
\\
\hline
\textbf{\textit{Use Case Name}}
&\textit{Login}
\\ \hline

\textbf{\textit{Actor}}
&\textit{Staff Infrastructure}, \textit{Manager Infrastructure}, \textit{Manager Human Resources} 
\\ \hline

\textbf{\textit{Brief Description}}
&\textit{Actor} yang ingin masuk ke dalam menu utama dan menggunakan seluruh fitur didalam aplikasi
\\ \hline

\textbf{\textit{Entry Conditions}}
&\textit{Actor} yang telah melakukan regitrasi akun dan akun sudah berstatus aktif
\\ \hline

\textbf{\textit{Flow of event}}
&1. \textit{Actor} membuka halaman \textit{login}
\newline
2. \textit{Actor} mengisi seluruh \textit{form input} login dan menekan tombol \textit{login}
\newline
3. Sistem melakukan pengecekkan terhadap \textit{credential} dari \textit{user}
\newline
4. Sistem regenerasi \textit{session} pada \textit{user}
\\ \hline

\textbf{\textit{Alternative scenario}}
&1. \textit{Actor} yang salah mengisi \textit{form input email} atau \textit{password} maka: Sistem menampilkan pesan error
\newline
2. \textit{Actor} yang sudah \textit{login} sebelumnya dan sistem akan mengecek \textit{session} dari \textit{user}, kemudian mengarahkan \textit{user} ke halaman utama
\\ \hline

\textbf{\textit{Exit Conditions}}
&Sistem menampilkan menu utama dengan fitur yang sesuai masing-masing \textit{position} dan \textit{department}
\\ \hline
\end{longtable}

%%%%%%%%%%%%%%%%%%%%%%%%%%%%%%%%%%%%%%%%%%%%%%%%%%%%%%%%%%%%%%%%%
\FloatBarrier
\begin{longtable}{ |p{6cm}|p{8cm}|}
\caption{\textit{Use Case Scenario Logout}}\label{tab:use-case-scenario-logout} 
\\
\hline
\textbf{\textit{Use Case Name}}
&\textit{Logout}
\\ \hline

\textbf{\textit{Actor}}
&\textit{Staff Infrastructure}, \textit{Manager Infrastructure}, \textit{Manager Human Resources} 
\\ \hline

\textbf{\textit{Brief Description}}
&\textit{Actor} yang telah melakukan \textit{login} dan ingin keluar dari sesi \textit{login}
\\ \hline

\textbf{\textit{Entry Conditions}}
&\textit{Actor} yang sudah melakukan \textit{login}
\\ \hline

\textbf{\textit{Flow of event}}
&1. \textit{Actor} menekan tombol yang bergambar \textit{profile user} pada \textit{navbar}
\newline
2. Sistem akan menampilkan beberapa pilihan, \textit{actor} menekan pilihan \textit{logout}
\newline
3. Sistem akan melakukan \textit{logout} dan menghapus \textit{session} dari \textit{user}
\\ \hline

\textbf{\textit{Alternative scenario}}
&-
\\ \hline

\textbf{\textit{Exit Conditions}}
&Sistem menampilkan menu \textit{login}
\\ \hline
\end{longtable}

%%%%%%%%%%%%%%%%%%%%%%%%%%%%%%%%%%%%%%%%%%%%%%%%%%%%%%%%%%%%%%%%%
\FloatBarrier
\begin{longtable}{ |p{6cm}|p{8cm}|}
\caption{\textit{Use Case Scenario Melakukan Clock In}}\label{tab:use-case-scenario-clock-in} 
\\
\hline
\textbf{\textit{Use Case Name}}
&\textit{Melakukan clock In}
\\ \hline

\textbf{\textit{Actor}}
&\textit{Staff Infrastructure}
\\ \hline

\textbf{\textit{Brief Description}}
&\textit{Actor} yang melakukan presensi masuk
\\ \hline

\textbf{\textit{Entry Conditions}}
&\textit{Actor} yang sudah melakukan \textit{login} dan belum melakukan \textit{clock in} sebelumnya
\\ \hline

\textbf{\textit{Flow of event}}
&1. \textit{Actor} menekan tombol \textit{presence} pada navbar
\newline
2. Sistem akan tampilan halaman \textit{presence}
\newline
3. \textit{Actor} menekan tombol \textit{clock in}
\newline
4. Sistem akan menampikan \textit{pop up modal} yang berisi tombol untuk pengambilan lokasi dan foto
\newline
5. \textit{Actor} menekan tombol \textit{get my location} untuk meng-\textit{input} lokasi \textit{user} secara otomatis melalui GPS dan \textit{Actor} menekan tombol \textit{take photo} untuk membuka kamera dan melakukan pengambilan foto
\newline
6. Mengisi \textit{form input} informasi presensi , kemudian menekan tombol \textit{clock in now}
\newline
7. Sistem akan menyimpan data presensi \textit{actor}
\\ \hline

\textbf{\textit{Alternative scenario}}
&\textit{Actor} yang sudah melakukan \textit{clock in} saat terakhir kali maka : tombol \textit{clock in} akan ter-\textit{disabled}
\\ \hline

\textbf{\textit{Exit Conditions}}
&Sistem selesai memproses data presensi masuk dan menampilkan \textit{pop up message} berhasil \textit{clock in}
\\ \hline
\end{longtable}


%%%%%%%%%%%%%%%%%%%%%%%%%%%%%%%%%%%%%%%%%%%%%%%%%%%%%%%%%%%%%%%%%
\FloatBarrier
\begin{longtable}{ |p{6cm}|p{8cm}|}
\caption{\textit{Use Case Scenario} Melakukan \textit{Clock Out}}\label{tab:use-case-scenario-clock-out} 
\\
\hline
\textbf{\textit{Use Case Name}}
&Melakukan \textit{clock Out}
\\ \hline

\textbf{\textit{Actor}}
&\textit{Staff Infrastructure}
\\ \hline

\textbf{\textit{Brief Description}}
&\textit{Actor} yang melakukan presensi keluar
\\ \hline

\textbf{\textit{Entry Conditions}}
&\textit{Actor} yang sudah melakukan \textit{login} dan sudah melakukan \textit{clock in} sebelumnya
\\ \hline

\textbf{\textit{Flow of event}}
&1. \textit{Actor} menekan tombol \textit{presence} pada navbar
\newline
2. Sistem akan tampilan halaman \textit{presence}
\newline
3. \textit{Actor} menekan tombol \textit{clock out}
\newline
4. Sistem akan menampikan \textit{pop up modal} yang berisi tombol untuk pengambilan lokasi dan foto
\newline
5. \textit{Actor} menekan tombol \textit{get my location} untuk meng-\textit{input} lokasi \textit{user} secara otomatis melalui GPS dan \textit{Actor} menekan tombol \textit{take photo} untuk membuka kamera dan melakukan pengambilan foto
\newline
6. Mengisi \textit{input form} informasi presensi, kemudian menekan tombol \textit{clock out now}
\newline
7. Sistem akan menyimpan data presensi \textit{actor}, dan sistem akan mengecek pada data presensi saat \textit{clock in} dan \textit{clock out} dan menyimpan data lembur jika presensi dilakukan diluar jam \textit{shift} kerja atau hari libur
\newline
8. Jika terdapat data lembur, sistem menyimpan data lembur
\\ \hline

\textbf{\textit{Alternative scenario}}
&\textit{Actor} yang belum melakukan \textit{clock in} saat terakhir kali maka : tombol \textit{clock out} akan ter-\textit{disabled}
\\ \hline

\textbf{\textit{Exit Conditions}}
&Sistem selesai memproses data presensi keluar dan data lembur, kemudian menampilkan \textit{pop up message} berhasil \textit{clock out}
\\ \hline
\end{longtable}


%%%%%%%%%%%%%%%%%%%%%%%%%%%%%%%%%%%%%%%%%%%%%%%%%%%%%%%%%%%%%%%%%
\FloatBarrier
\begin{longtable}{ |p{6cm}|p{8cm}|}
\caption{\textit{Use Case Scenario} Melakukan Pengajuan Lembur}\label{tab:use-case-scenario-melakukan-pengajuan-lembur} 
\\
\hline
\textbf{\textit{Use Case Name}}
&Melakukan pengajuan lembur
\\ \hline

\textbf{\textit{Actor}}
&\textit{Staff Infrastructure}
\\ \hline

\textbf{\textit{Brief Description}}
&\textit{Actor} yang melakukan pengajuan lembur
\\ \hline

\textbf{\textit{Entry Conditions}}
&\textit{Actor} yang sudah melakukan \textit{login} dan ingin mengajukan lembur karena lupa presensi saat jam lembur
\\ \hline

\textbf{\textit{Flow of event}}
&1. \textit{Actor} menekan tombol \textit{request} pada navbar
\newline
2. Sistem akan menampilkan halaman \textit{overtime correction request}
\newline
3. \textit{Actor} mengisi seluruh \textit{form input} pada halaman \textit{overtime correction request} dan menekan tombol submit
\newline
4. Sistem akan menyimpan data lembur yang telah diisi oleh \textit{actor}
\\ \hline

\textbf{\textit{Alternative scenario}}
&\textit{Actor} yang tidak mengisi seluruh \textit{input form} maka : Sistem akan memberikan pesan error
\\ \hline

\textbf{\textit{Exit Conditions}}
&Sistem selesai memproses data lembur dan menampilkan \textit{pop up message} berhasil dan mengembalikan tampilan ke halaman \textit{overtime correction request}
\\ \hline
\end{longtable}

%%%%%%%%%%%%%%%%%%%%%%%%%%%%%%%%%%%%%%%%%%%%%%%%%%%%%%%%%%%%%%%%%
\FloatBarrier
\begin{longtable}{ |p{6cm}|p{8cm}|}
\caption{\textit{Use Case Scenario} Melihat Jadwal \textit{Shift} Kerja}\label{tab:use-case-scenario-melihat-jadwal-shift-kerja} 
\\
\hline
\textbf{\textit{Use Case Name}}
&Melihat jadwal \textit{shift} kerja
\\ \hline

\textbf{\textit{Actor}}
&\textit{Staff Infrastructure}
\\ \hline

\textbf{\textit{Brief Description}}
&\textit{Actor} yang melihat jadwal \textit{shift} kerja
\\ \hline

\textbf{\textit{Entry Conditions}}
&\textit{Actor} yang sudah melakukan \textit{login} dan ingin melihat jadwal \textit{shift} kerja
\\ \hline

\textbf{\textit{Flow of event}}
&1. \textit{Actor} menekan tombol \textit{presence} pada \textit{navbar}
\newline
2. Sistem akan menampilkan halaman \textit{presence}
\newline
3. \textit{Actor} menekan tombol \textit{shift schedule}
\newline
4. Sistem akan menampilkan jadwal kerja \textit{actor}
\\ \hline

\textbf{\textit{Alternative scenario}}
&\textit{Actor} yang tidak memiliki jadwal kerja maka tidak akan ada data yang ditampilkan
\\ \hline

\textbf{\textit{Exit Conditions}}
&\textit{Actor} menutup \textit{modal} data \textit{shift} dan sistem menampilkan kembali halaman \textit{attendance}
\\ \hline
\end{longtable}

%%%%%%%%%%%%%%%%%%%%%%%%%%%%%%%%%%%%%%%%%%%%%%%%%%%%%%%%%%%%%%%%%
\FloatBarrier
\begin{longtable}{ |p{6cm}|p{8cm}|}
\caption{\textit{Use Case Scenario} Melakukan \textit{Reject} Pada Data Lembur}\label{tab:use-case-scenario-melakukan-approval} 
\\
\hline
\textbf{\textit{Use Case Name}}
&Melakukan \textit{reject} pada data lembur
\\ \hline

\textbf{\textit{Actor}}
&\textit{Manager Infrastructure}
\\ \hline

\textbf{\textit{Brief Description}}
&\textit{Actor} yang melakukan \textit{reject} pada data lembur
\\ \hline

\textbf{\textit{Entry Conditions}}
&\textit{Actor} yang sudah melakukan \textit{login} dan ingin melakukan penolakkan atau \textit{reject} pada data lembur yang masuk
\\ \hline

\textbf{\textit{Flow of event}}
&1. \textit{Actor} menekan tombol \textit{reject} pada \textit{navbar}
\newline
2. Sistem akan menampilkan \textit{list} data \textit{lembur} yang tidak pernah di-\textit{reject}
\newline
3. \textit{Actor} melakukan \textit{reject} pada data yang dipilih atau mencetang \textit{checkbox} pada data lembur dan menekan tombol \textit{reject selected}
\newline
5. \textit{Actor} melakukan konfirmasi perubahan
\newline
6. Sistem menyimpan perubahan data
\\ \hline

\textbf{\textit{Alternative scenario}}
&-
\\ \hline

\textbf{\textit{Exit Conditions}}
&Sistem selesai memproses seluruh perubahan data lembur dan kembali menampilkan halaman \textit{reject}
\\ \hline
\end{longtable}

%%%%%%%%%%%%%%%%%%%%%%%%%%%%%%%%%%%%%%%%%%%%%%%%%%%%%%%%%%%%%%%%%
\FloatBarrier
\begin{longtable}{ |p{6cm}|p{8cm}|}
\caption{\textit{Use Case Scenario} Melakukan \textit{Export} Data Lembur ke \textit{Excel}}\label{tab:use-case-scenario-melakukan-export-data-lembur} 
\\
\hline
\textbf{\textit{Use Case Name}}
&Melakukan \textit{Export} Data Lembur ke \textit{Excel}
\\ \hline

\textbf{\textit{Actor}}
&\textit{Manager Human Resources, Manager Infrastructure}
\\ \hline

\textbf{\textit{Brief Description}}
&\textit{Actor} yang melakukan \textit{export} data lembur menjadi format \textit{Excel}
\\ \hline

\textbf{\textit{Entry Conditions}}
&\textit{Actor} yang sudah melakukan \textit{login} dan ingin melakukan \textit{export} data lembur menjadi format \textit{Excel}
\\ \hline

\textbf{\textit{Flow of event}}
&1. \textit{Actor} menekan tombol \textit{overtime history} pada \textit{navbar}
\newline
2. Sistem akan menampilkan \textit{list} data riwayat lembur
\newline
3. \textit{Actor} melakukan \textit{filter} nama \textit{staff}, tanggal atau status
\newline
4. Sistem menampilkan data lembur yang telah di-\textit{filter}
\newline
5. \textit{Actor} menekan tombol \textit{Export}
\newline
6. Sistem akan memproses data lembur dan memberikan konfirmasi untuk men-\textit{download} \textit{file} dalam bentu \textit{Excel}
\newline
7. \textit{Actor} mengunduh \textit{file Excel}
\newline
\\ \hline

\textbf{\textit{Alternative scenario}}
&Jika tidak ada data lembur atau data yang di-\textit{filter} tidak ditemukan, maka sistem tidak akan memproses permintaan \textit{export}
\\ \hline

\textbf{\textit{Exit Conditions}}
&\textit{Actor} menekan konfirmasi saat pengunduhan \textit{file} dan sistem akan mengembalikan tampilan ke \textit{overtime history}
\\ \hline
\end{longtable}

%%%%%%%%%%%%%%%%%%%%%%%%%%%%%%%%%%%%%%%%%%%%%%%%%%%%%%%%%%%%%%%%%
\FloatBarrier
\begin{longtable}{ |p{6cm}|p{8cm}|}
\caption{\textit{Use Case Scenario} Melakukan \textit{Export} Data Presensi ke \textit{Excel}}\label{tab:use-case-scenario-melakukan-export-data-presensi} 
\\
\hline
\textbf{\textit{Use Case Name}}
&Melakukan \textit{Export} Data Presensi ke \textit{Excel}
\\ \hline

\textbf{\textit{Actor}}
&\textit{Manager Human Resources, Manager Infrastructure}
\\ \hline

\textbf{\textit{Brief Description}}
&\textit{Actor} yang melakukan \textit{export} data presensi menjadi format \textit{Excel}
\\ \hline

\textbf{\textit{Entry Conditions}}
&\textit{Actor} yang sudah melakukan \textit{login} dan ingin melakukan \textit{export} data presensi menjadi format \textit{Excel}
\\ \hline

\textbf{\textit{Flow of event}}
&1. \textit{Actor} menekan tombol \textit{attendance history} pada \textit{navbar}
\newline
2. Sistem akan menampilkan \textit{list} data riwayat presensi
\newline
3. \textit{Actor} melakukan \textit{filter} nama \textit{staff}, tanggal atau status
\newline
4. Sistem menampilkan data presensi yang telah di-\textit{filter}
\newline
5. \textit{Actor} menekan tombol \textit{Export}
\newline
6. Sistem akan memproses data presensi dan memberikan konfirmasi untuk men-\textit{download} \textit{file} dalam bentu \textit{Excel}
\newline
7. \textit{Actor} mengunduh \textit{file Excel}
\newline
\\ \hline

\textbf{\textit{Alternative scenario}}
&Jika tidak ada data presensi atau data yang di-\textit{filter} tidak ditemukan, maka sistem tidak akan memproses permintaan \textit{export}
\\ \hline

\textbf{\textit{Exit Conditions}}
&\textit{Actor} menekan konfirmasi saat pengunduhan \textit{file} dan sistem akan mengembalikan tampilan ke \textit{attendance history}
\\ \hline
\end{longtable}

%%%%%%%%%%%%%%%%%%%%%%%%%%%%%%%%%%%%%%%%%%%%%%%%%%%%%%%%%%%%%%%%%
\FloatBarrier
\begin{longtable}{ |p{6cm}|p{8cm}|}
\caption{\textit{Use Case Scenario} Melihat Riwayat Lembur}\label{tab:use-case-scenario-melihat-riwayat-lembur} 
\\
\hline
\textbf{\textit{Use Case Name}}
&Melihat riwayat lembur
\\ \hline

\textbf{\textit{Actor}}
&\textit{Staff Infrastructure}, \textit{Manager Infrastructure}, \textit{Manager Human Resources}
\\ \hline

\textbf{\textit{Brief Description}}
&\textit{Actor} yang melihat riwayat lembur
\\ \hline

\textbf{\textit{Entry Conditions}}
&\textit{Actor} yang sudah melakukan \textit{login} dan ingin melihat riwayat lembur
\\ \hline

\textbf{\textit{Flow of event}}
&1. \textit{Actor} menekan tombol \textit{overtime history} pada \textit{navbar}
\newline
2. Sistem akan menampilkan \textit{list} data riwayat lembur
\newline
3. \textit{Actor} melakukan \textit{filter} tanggal dan status pada data lembur untuk mencari spesifik data lembur
\newline
4. \textit{Actor} menekan tombol \textit{details} pada salah satu data lembur
\newline
5. Sistem menampilkan halaman \textit{details} dan informasi lengkap terkait data lembur yang dipilih
\newline
\\ \hline

\textbf{\textit{Alternative scenario}}
&Jika tidak ada data lembur atau data yang di-\textit{filter} tidak ditemukan, maka sistem tidak akan menampilkan data
\\ \hline

\textbf{\textit{Exit Conditions}}
&\textit{Actor} menekan tombol kembali dan sistem akan menampilkan kembali halaman \textit{overtime history}
\\ \hline
\end{longtable}

%%%%%%%%%%%%%%%%%%%%%%%%%%%%%%%%%%%%%%%%%%%%%%%%%%%%%%%%%%%%%%%%%
\FloatBarrier
\begin{longtable}{ |p{6cm}|p{8cm}|}
\caption{\textit{Use Case Scenario} Melihat Riwayat Presensi}\label{tab:use-case-scenario-melihat-riwayat-presensi} 
\\
\hline
\textbf{\textit{Use Case Name}}
&Melihat riwayat presensi
\\ \hline

\textbf{\textit{Actor}}
&\textit{Staff Infrastructure}, \textit{Manager Infrastructure}, \textit{Manager Human Resources}
\\ \hline

\textbf{\textit{Brief Description}}
&\textit{Actor} yang melihat riwayat presensi
\\ \hline

\textbf{\textit{Entry Conditions}}
&\textit{Actor} yang sudah melakukan \textit{login} dan ingin melihat riwayat presensi
\\ \hline

\textbf{\textit{Flow of event}}
&1. \textit{Actor} menekan tombol \textit{attendance history} pada \textit{navbar}
\newline
2. Sistem akan menampilkan \textit{list} data riwayat presensi
\newline
3. \textit{Actor} melakukan \textit{filter} tanggal pada data presensi untuk mencari spesifik data presensi
\newline
4. \textit{Actor} menekan tombol \textit{details} pada salah satu data presensi
\newline
5. Sistem menampilkan halaman \textit{details} dan informasi lengkap terkait data presensi yang dipilih
\newline
\\ \hline

\textbf{\textit{Alternative scenario}}
&Jika tidak ada data presensi atau data yang di-\textit{filter} tidak ditemukan, maka sistem tidak akan menampilkan data
\\ \hline

\textbf{\textit{Exit Conditions}}
&\textit{Actor} menekan tombol kembali dan sistem akan menampilkan kembali halaman \textit{attendance history}
\\ \hline
\end{longtable}

%%%%%%%%%%%%%%%%%%%%%%%%%%%%%%%%%%%%%%%%%%%%%%%%%%%%%%%%%%%%%%%%%
\FloatBarrier
\begin{longtable}{ |p{6cm}|p{8cm}|}
\caption{\textit{Use Case Scenario} Mengelola Jenis \textit{Shift}}\label{tab:use-case-scenario-mengelola-jenis-shift} 
\\
\hline
\textbf{\textit{Use Case Name}}
&Mengelola jenis \textit{shift}
\\ \hline

\textbf{\textit{Actor}}
&\textit{Manager Human Resources}, \textit{Manager Infrastructure}
\\ \hline

\textbf{\textit{Brief Description}}
&\textit{Actor} yang mengelola jenis \textit{shift}
\\ \hline

\textbf{\textit{Entry Conditions}}
&\textit{Actor} yang sudah melakukan \textit{login} dan ingin mengelola jenis \textit{shift}
\\ \hline

\textbf{\textit{Flow of event}}
&1. \textit{Actor} menekan tombol \textit{shift} pada \textit{navbar}
\newline
2. Sistem akan menampilkan halaman \textit{shift} beserta data jenis \textit{shift} kerja
\newline
3. \textit{Actor} mengelola (\textit{create, update, delete}) data jenis \textit{shift}
    \newline
4. Sistem menampilkan konfirmasi perubahan data
\newline
5. \textit{Actor} melakukan konfirmasi perubahan
\newline
6. Sistem menyimpan perubahan data
\newline
\\ \hline

\textbf{\textit{Alternative scenario}}
&\textit{Actor} yang tidak melakukan konfirmasi perubahan data maka sistem tidak akan menyimpan perubahan data
\\ \hline

\textbf{\textit{Exit Conditions}}
&Sistem selesai memproses seluruh perubahan data jenis \textit{shift} dan kembali menampilkan halaman \textit{shift}
\\ \hline
\end{longtable}

%%%%%%%%%%%%%%%%%%%%%%%%%%%%%%%%%%%%%%%%%%%%%%%%%%%%%%%%%%%%%%%%%
\FloatBarrier
\begin{longtable}{ |p{6cm}|p{8cm}|}
\caption{\textit{Use Case Scenario} Mengelola Hari \textit{Shift}}\label{tab:use-case-scenario-mengelola-hari-shift} 
\\
\hline
\textbf{\textit{Use Case Name}}
&Mengelola hari \textit{shift}
\\ \hline

\textbf{\textit{Actor}}
&\textit{Manager Human Resources}, \textit{Manager Infrastructure}
\\ \hline

\textbf{\textit{Brief Description}}
&\textit{Actor} yang mengelola hari \textit{shift}
\\ \hline

\textbf{\textit{Entry Conditions}}
&\textit{Actor} yang sudah melakukan \textit{login} dan ingin mengelola hari \textit{shift}
\\ \hline

\textbf{\textit{Flow of event}}
&1. \textit{Actor} menekan \textit{edit} pada salah satu jenis \textit{shift} melalui halaman \textit{shift}
\newline
2. Sistem akan menampilkan \textit{modal} hari \textit{shift} beserta data hari \textit{shift} kerja
\newline
3. \textit{Actor} mengelola (\textit{create, update, delete}) data hari \textit{shift}
    \newline
4. Sistem menampilkan konfirmasi perubahan data
\newline
5. \textit{Actor} melakukan konfirmasi perubahan
\newline
6. Sistem menyimpan perubahan data
\newline
\\ \hline

\textbf{\textit{Alternative scenario}}
&\textit{Actor} yang tidak melakukan konfirmasi perubahan data maka sistem tidak akan menyimpan perubahan data
\\ \hline

\textbf{\textit{Exit Conditions}}
&Sistem selesai memproses seluruh perubahan data hari \textit{shift} dan kembali menampilkan halaman \textit{shift}
\\ \hline
\end{longtable}

%%%%%%%%%%%%%%%%%%%%%%%%%%%%%%%%%%%%%%%%%%%%%%%%%%%%%%%%%%%%%%%%%
\FloatBarrier
\begin{longtable}{ |p{6cm}|p{8cm}|}
\caption{\textit{Use Case Scenario} Mengelola Penjadwalan \textit{Shift}}\label{tab:use-case-scenario-mengelola-penjadwalan-shift}
\\
\hline
\textbf{\textit{Use Case Name}}
&Mengelola penjadwalan \textit{shift}
\\ \hline

\textbf{\textit{Actor}}
&\textit{Manager Human Resources}, \textit{Manager Infrastructure}
\\ \hline

\textbf{\textit{Brief Description}}
&\textit{Actor} yang mengelola penjadwalan \textit{shift}
\\ \hline

\textbf{\textit{Entry Conditions}}
&\textit{Actor} yang sudah melakukan \textit{login} dan ingin mengelola penjadwalan \textit{shift}
\\ \hline

\textbf{\textit{Flow of event}}
&1. \textit{Actor} menekan tombol \textit{shift scheduling} pada \textit{navbar}
\newline
2. Sistem akan menampilkan halaman penjadwalan \textit{shift} beserta data penjadwalan \textit{shift} kerja
\newline
3. \textit{Actor} mengelola (\textit{create, update, delete, import}) data penjadwalan \textit{shift}
    \newline
4. Sistem menampilkan konfirmasi perubahan data
\newline
5. \textit{Actor} melakukan konfirmasi perubahan
\newline
6. Sistem menyimpan perubahan data
\newline
\\ \hline

\textbf{\textit{Alternative scenario}}
&\textit{Actor} yang tidak melakukan konfirmasi perubahan data maka sistem tidak akan menyimpan perubahan data
\\ \hline

\textbf{\textit{Exit Conditions}}
&Sistem selesai memproses seluruh perubahan data penjadwalan \textit{shift} dan kembali menampilkan halaman penjadwalan \textit{shift}
\\ \hline
\end{longtable}


%%%%%%%%%%%%%%%%%%%%%%%%%%%%%%%%%%%%%%%%%%%%%%%%%%%%%%%%%%%%%%%%%
\FloatBarrier
\begin{longtable}{ |p{6cm}|p{8cm}|}
\caption{\textit{Use Case Scenario} Mengelola Akun \textit{User}}\label{tab:use-case-scenario-mengelola-akun-user} 
\\
\hline
\textbf{\textit{Use Case Name}}
&Mengelola akun \textit{user}
\\ \hline

\textbf{\textit{Actor}}
&\textit{Manager Human Resources}, \textit{Manager Infrastructure}
\\ \hline

\textbf{\textit{Brief Description}}
&\textit{Actor} yang mengelola akun \textit{user}
\\ \hline

\textbf{\textit{Entry Conditions}}
&\textit{Actor} yang sudah melakukan \textit{login} dan ingin mengelola akun \textit{user}
\\ \hline

\textbf{\textit{Flow of event}}
&1. \textit{Actor} menekan tombol \textit{user} pada \textit{navbar}
\newline
2. Sistem akan menampilkan halaman \textit{user} beserta data semua \textit{user}
\newline
3. \textit{Actor} mengelola \textit{update, delete} data akun \textit{user}
\newline
4. Sistem menampilkan konfirmasi perubahan data
\newline
5. \textit{Actor} melakukan konfirmasi perubahan
\newline
6. Sistem menyimpan perubahan data
\\ \hline

\textbf{\textit{Alternative scenario}}
&\textit{Actor} yang tidak melakukan konfirmasi perubahan data maka sistem tidak akan menyimpan perubahan data
\\ \hline

\textbf{\textit{Exit Conditions}}
&Sistem selesai memproses seluruh perubahan data \textit{user} dan kembali menampilkan halaman \textit{user}
\\ \hline
\end{longtable}

%%%%%%%%%%%%%%%%%%%%%%%%%%%%%%%%%%%%%%%%%%%%%%%%%%%%%%%%%%%%%%%%%
\FloatBarrier
\begin{longtable}{ |p{6cm}|p{8cm}|}
\caption{\textit{Use Case Scenario} Mengelola \textit{Department}}\label{tab:use-case-scenario-mengelola-department} 
\\
\hline
\textbf{\textit{Use Case Name}}
&Mengelola \textit{department}
\\ \hline

\textbf{\textit{Actor}}
&\textit{Manager Human Resources}, \textit{Manager Infrastructure}
\\ \hline

\textbf{\textit{Brief Description}}
&\textit{Actor} yang mengelola \textit{department}
\\ \hline

\textbf{\textit{Entry Conditions}}
&\textit{Actor} yang sudah melakukan \textit{login} dan ingin mengelola \textit{department}
\\ \hline

\textbf{\textit{Flow of event}}
&1. \textit{Actor} menekan tombol \textit{department} pada \textit{navbar}
\newline
2. Sistem akan menampilkan halaman \textit{department} beserta data semua \textit{department}
\newline
3. \textit{Actor} mengelola \textit{update, delete} data akun \textit{department}
\newline
4. Sistem menampilkan konfirmasi perubahan data
\newline
5. \textit{Actor} melakukan konfirmasi perubahan
\newline
6. Sistem menyimpan perubahan data
\\ \hline

\textbf{\textit{Alternative scenario}}
&\textit{Actor} yang tidak melakukan konfirmasi perubahan data maka sistem tidak akan menyimpan perubahan data
\\ \hline

\textbf{\textit{Exit Conditions}}
&Sistem selesai memproses seluruh perubahan data \textit{department} dan kembali menampilkan halaman \textit{department}
\\ \hline
\end{longtable}

%%%%%%%%%%%%%%%%%%%%%%%%%%%%%%%%%%%%%%%%%%%%%%%%%%%%%%%%%%%%%%%%%
\FloatBarrier
\begin{longtable}{ |p{6cm}|p{8cm}|}
\caption{\textit{Use Case Scenario} Mengelola Jenis Libur Perusahaan}\label{tab:use-case-scenario-mengelola-jenis-libur-perusahaan} 
\\
\hline
\textbf{\textit{Use Case Name}}
&Mengelola jenis libur perusahaan
\\ \hline

\textbf{\textit{Actor}}
&\textit{Manager Human Resources}, \textit{Manager Infrastructure}
\\ \hline

\textbf{\textit{Brief Description}}
&\textit{Actor} yang mengelola jenis libur perusahaan
\\ \hline

\textbf{\textit{Entry Conditions}}
&\textit{Actor} yang sudah melakukan \textit{login} dan ingin mengelola jenis libur perusahaan
\\ \hline

\textbf{\textit{Flow of event}}
&1. \textit{Actor} menekan tombol libur pada \textit{navbar}
\newline
2. Sistem akan menampilkan halaman libur beserta seluruh data jenis libur perusahaan
\newline
3. \textit{Actor} mengelola (\textit{create, update, delete}) data jenis libur perusahaan
\newline
4. Sistem menampilkan konfirmasi perubahan data
\newline
5. \textit{Actor} melakukan konfirmasi perubahan
\newline
6. Sistem menyimpan perubahan data
\\ \hline

\textbf{\textit{Alternative scenario}}
&1. \textit{Actor} yang tidak melakukan konfirmasi perubahan data maka sistem tidak akan menyimpan perubahan data
\newline
2. Jika tidak ada data hari libur yang ditemukan, maka sistem tidak akan menampilkan hasil
\\ \hline

\textbf{\textit{Exit Conditions}}
&Sistem selesai memproses seluruh perubahan data jenis libur perusahaan dan kembali menampilkan halaman jenis libur
\\ \hline
\end{longtable}

%%%%%%%%%%%%%%%%%%%%%%%%%%%%%%%%%%%%%%%%%%%%%%%%%%%%%%%%%%%%%%%%%
\FloatBarrier
\begin{longtable}{ |p{6cm}|p{8cm}|}
\caption{\textit{Use Case Scenario} Mengelola Data Hari Libur Perusahaan}\label{tab:use-case-scenario-mengelola-data-hari-libur-perusahaan} 
\\
\hline
\textbf{\textit{Use Case Name}}
&Mengelola data hari libur perusahaan
\\ \hline

\textbf{\textit{Actor}}
&\textit{Manager Human Resources}, \textit{Manager Infrastructure}
\\ \hline

\textbf{\textit{Brief Description}}
&\textit{Actor} yang mengelola data hari libur perusahaan
\\ \hline

\textbf{\textit{Entry Conditions}}
&\textit{Actor} yang sudah melakukan \textit{login} dan ingin mengelola data hari libur perusahaan
\\ \hline

\textbf{\textit{Flow of event}}
&1. \textit{Actor} menekan tombol edit pada jenis libur perusahaan
\newline
2. Sistem akan menampilkan halaman hari libur beserta seluruh data hari libur perusahaan
\newline
3. \textit{Actor} mengelola (\textit{create, update, delete}) data hari libur perusahaan atau \textit{import} data hari libur nasional melalui tombol yang terhubung dengan API
\newline
4. Sistem menampilkan konfirmasi perubahan data
\newline
5. \textit{Actor} melakukan konfirmasi perubahan
\newline
6. Sistem menyimpan perubahan data
\\ \hline

\textbf{\textit{Alternative scenario}}
&1. \textit{Actor} yang tidak melakukan konfirmasi perubahan data maka sistem tidak akan menyimpan perubahan data
\newline
2. Jika tidak ada data hari libur yang ditemukan, maka sistem tidak akan menampilkan hasil
\\ \hline

\textbf{\textit{Exit Conditions}}
&Sistem selesai memproses seluruh perubahan data hari libur perusahaan dan kembali menampilkan halaman hari libur
\\ \hline
\end{longtable}

